\documentclass[letterpaper, 10pt, conference]{ieeeconf}  % 
\IEEEoverridecommandlockouts                              % This command is only needed if 
                                                          % you want to use the \thanks command
\overrideIEEEmargins    
\usepackage{cite}

\usepackage{algorithmic}

\usepackage{textcomp}

\usepackage{comment}

%\usepackage{graphicx}      % include this line if your document contains figures
%\usepackage{natbib}        % required for bibliography
%===============================================================================
% Needed to meet printer requirements.
 \usepackage{latexsym,amssymb,amsfonts}
\usepackage[applemac]{inputenc}
\usepackage{rotating}
\usepackage{cite}
%\usepackage{graphicx}
\usepackage{epstopdf}
\usepackage{subfigure}
%\usepackage{color}
%\usepackage{xcolor}
%
%\usepackage{color}
\usepackage[dvipsnames]{xcolor}
%
\usepackage[normalem]{ulem}
\usepackage{isolatin1}
%\usepackage{caption}
%\usepackage{array,multirow,makecell}
%\usepackage{placeins}
%\usepackage{stackengine}
%\usepackage{hyperref}
\usepackage{fancybox}
\usepackage{colortbl}
\usepackage{float}
\usepackage{comment}
\usepackage{mathdots}
%\usepackage{lmodern}
%%%%%%%%%%%%%%%%%%%%%%%%%%%%%%%%%%%%%%%%%%%%%%%%%%%%%%%%%
\usepackage{dblfloatfix}
\usepackage{kantlipsum}
\usepackage{dblfloatfix}    % To enable figures at the bottom of page
\usepackage{graphicx} % [demo] option for empty figure
\usepackage{placeins}
\usepackage{marginnote}
\usepackage{url}
\usepackage{hyperref}
\usepackage{fancybox}
%\usepackage{empheq}
\usepackage[overload]{empheq}

\usepackage{bm}

%%%%%%%%%%%%%%%%%%%%%%%%%%%%%%%%%%%%%%%%%%%%%%%%%%%%%%%%%
\usepackage{tikz}
\usetikzlibrary{circuits.ee.IEC}
\usepackage[european resistor, european voltage, european current]{circuitikz}
\usetikzlibrary{arrows,shapes,positioning}
\usetikzlibrary{decorations.markings,decorations.pathmorphing,
decorations.pathreplacing}
\usetikzlibrary{calc,patterns,shapes.geometric}
\usetikzlibrary{shapes.multipart,positioning,patterns,backgrounds}
\tikzstyle{titlebox}=[rectangle,inner sep=10pt,inner ysep=10pt,draw]%

%%%%%%%%%%%%%%%%%%%%%%%%%%%%%%%%%%%%%%%%%%%%%%%%%%%%%%%%%%%%%

\def\blacksquare{\hbox{\vrule width 4pt height 4pt depth 0pt}}
\def\square{\hbox{\vrule\barox{\hrule\phantom{o}\hrule}\vrule}}
\def\qed{\hfill \square}
\def\qedp{\rightline{$\blacksquare$}}
\def\EXPT{\mathop{{\rm E}}}
\def\MIN{\mathop{{\rm min}}}
\def\MAX{\mathop{{\rm max}}}
\def\LIM{\mathop{{\longrightarrow}}}
\def\UNION{\mathop{\bigcup}}
\def\INTERSECT{\mathop{\bigcap}}
\def\ARGMIN{\mathop{{\rm argmin}}}
\def\ARGMAX{\mathop{{\rm argmax}}}
\def\TOUND{\mathop{\longrightarrow}}
\def\col#1{\mathop{{\rm col}\,\left(\,#1\,\right)}}
\def\eqref#1{(\ref{#1})}
\def\rank{\mathop{{\rm rank}}}

\newcommand{\widebox}[2][0.5em]{\fbox{\hspace{#1}$\displaystyle #2$\hspace{#1}}}


%\newcommand*\rank{\rm rank}
\def\rank{\mathop{{\rm rank}}}

%%%
\newtheorem{prb}{\it Problem}%[section]
\newtheorem{theorem}{\it Theorem}%[section]
\newtheorem{lemma}{\it Lemma}%[section]
\newtheorem{definition}{\it Definition}
\newtheorem{remark}{\it Remark}%[section]
\newtheorem{assumption}{\it Assumption}%[section]
\newtheorem{example}{\it Example}%[section]b
\newtheorem{annex}{\it Annex}%[section]b
\newtheorem{prop}{\it Proposition}%[section]b








\def\BibTeX{{\rm B\kern-.05em{\sc i\kern-.025em b}\kern-.08em
    T\kern-.1667em\lower.7ex\hbox{E}\kern-.125emX}}

\begin{document}


\title{\bf Resilience in Distributed
 Consensus based on Trust for Conneted Vehicle }
\author{Q.H. Nguyen$^{1,2}$, Shengya Meng $^{1}$, H. Rafaralahy $^{1}$, M. Haddad $^{2}$, A. Zemouche $^{1}$
%
\thanks{This work was supported by SEGULA Engineering. H. Rafaralahy and A. Zemouche would like to acknowledge the ANR agency for its partial financial support under the project ArtISMo ANR-20-CE48-0015.
}
%
\thanks{$^{1}$ Universit\'e de Lorraine, CRAN  CNRS UMR 7039, 54400 Cosnes et Romain, France~(e-mail: \protect\url{ali.zemouche@univ-lorraine.fr}).}
%
\thanks{$^{2}$ SEGULA Engineering, 19 rue d'Arras, 92000 Nanterre, France~(e-mail: \protect\url{madjid.haddad@segula.fr}).}
%
\thanks{All relevant CARLA resources have been made available in \protect\href{https://github.com/kslhuy/CACC_UIO_simulink}{GitHub}.}
}
\maketitle

\begin{abstract}
We propose a distributed driving state observer inspired by the leader-follower consensus technique. This observer reconstructs the global driving state of a vehicle platoon, which includes the positions, velocities, and accelerations of all vehicles. Additionally, we establish the necessary and sufficient conditions to ensure the stability of the estimation error dynamics. Our work addresses the challenge of designing a resilient distributed observer in the presence of cyber-attacks by calculating stochastic trust values among agents. This allows us to design resilient weights that can exclude malicious agents from the estimation process, thereby preserving the integrity of the good estimations.
\end{abstract}

\begin{keywords}
Distributed observers, secure estimation, lyapunov stability, cyber-attack.
\end{keywords}

\section{Introduction}

 The objective of global driving state
 estimation is to design an estimation algorithm for each vehicle
 to determine the positions, velocities, and accelerations of all
 vehicles in the platoon using only locally available information

 As a platoon is a typical networked control system (NCS),
 there has been a recent surge of interest in distributed driving
 state estimation

\subsection{Related Work}
There has been significant research recently on
 cyber-security of CACC. In this section, we summarize
 the research in different categories of security solutions for
 CACC and highlight their limitations. Additionally, several
 excellent surveys are available that provide a comprehensive
 treatment of attacks on CAVs engaging in cooperative driving
 applications [14], [29], [30], [36].
 \paragraph*{(1) Offline Detection Techniques}
  : Biroon et al. [7] propose
 a diagnostic scheme based on a partial differential equation
 observer model to detect an adversary injecting ghost vehi
cles into the platoon. Keijzer and Ferrari [20] propose a
 sliding-mode observer (SMO) approach for detecting attacks
 on CACC-based platoon compromising V2V and local sensors
 simultaneously. Jagielski et al. [16] present a discussion on
 detection of communication and sensor attacks compromis
ing CACC. Alotibi and Abdelhakim [3] propose a detection
 technique to capture falsified leader communication in a pla
toon. Mokari et al. [26] propose a detection and estimation
 approach against Denial of Service (DoS) attacks on platoons.
 Sajjad et al. [34] propose an adversarial-aware CACC control
 scheme that utilizes only local sensor information obviating
 the need for inter-vehicle communication. This category of
 techniques target detection of anomalous samples in historic
 data after application engagement and therefore, ineffective
 for real-time resiliency.
 \paragraph*{2) Consensus-Based Approaches:} Lu et al. [23] present
 attack-resilient sensor fusion approach for platoons by uti
lizing the spatial information provided by other participating
 CAVs thereby achieving more accurate estimation. Yang and
 Lv [42] propose adding redundant sensor systems to obtain
 robust estimates of the physical parameters susceptible to
 security attacks. Petrillo et al. [31] present a secure adap
tive CACC controller capable of evicting malicious CAVs
 in the platoon and mitigating the effects of network-induced
 perturbations. Kamel et al. [19] propose a simulation frame
work for misbehavior detection in vehicular platoons through
 plausibility and consistency checks of the reported messages.
 Garlichs et al. [13] propose a trust model to identify bad
 actors and prevent safety hazards in dynamically formed
 platoon systems.

 Obviously, approaches relying on majority voting cannot
 apply to 2-vehicle CACC. These are also ineffective in the
 absence of a fixed trusted source of information.
 \paragraph*{ 3) Mitigation Techniques:} Khanapuri et al. [21] propose
 a detection and mitigation approach for defending pla
toon systems against attacks on V2V communication. They
 utilize local sensor information to mitigate V2V attacks.
 Wolf et al. [39] propose various general mitigation strate
gies against data injection attacks and perception channel
 faults. According to one of the detection-oriented approaches
 suggested here, the CACC controller gradually degrades
 to ACC in response to the suspiciousness index of the
 received inputs. Iorio et al. [15] present a correlation-based
 anomaly detection technique for capturing injection attacks
 on CACC. Jin et al. [17] develop an adaptive CACC con
troller that can defend against sensor and actuator attacks.
 Sun et al. [35] present a detection and mitigation approach
 for V2V attacks on CACC based on deep-learning classifiers.
 Kalogiannis et al. [18], propose a misbehavior detection and
 mitigation technique for platoons based on Guassian Mixture
 Model, under a collaborative adversary with an internal and
 external attack agent acting together to subvert the platoon
 system. Mousavinejad et al. [27], propose a distributed detec
tion and recovery method based on state prediction and state
 estimation to identify malicious platoon messages. Each of
 these approaches generally targets a specific attack category
 (e.g., DoS or injection). It is difficult to integrate individual
 resiliency solutions into vehicular electronics to address each
 category of attack instances. Furthermore, such solutions can
not provide resiliency against attacks that are discovered after
 deployment. Previous work by Boddupalli et al. [8] addresses
 some of these limitations. They develop an ML-based real
time resiliency architecture for 2-vehicle CACC to capture
 and mitigate anomalies in the untrusted V2V channel reporting
 preceding vehicle acceleration values. A unique feature of this
 work is that it provides assured resiliency against attacks on
 V2X irrespective of the attack mechanism; consequently, the
 protection also extends to unknown attacks. This work is most
 closely related to our work in this paper and provides an initial
 roadmap for developing resiliency architecture for CACC
 that we extend and consolidate with RECAP. However, the
 adversary considered in this previous work could only corrupt
 the acceleration channel; all other channels (e.g., preceding
 velocity and gap between the ego and preceding vehicles) were
 assumed to report ground truth. Their mitigation critically
 depends on this fact to compute an alternative response in
 case an anomaly is detected. Consequently, such an approach
 is precluded under the multi-channel adversary considered for
 RECAP where no channel can be assumed trusted

\paragraph{Contributions}
Contributions in the Context of Previous Work: Instead, the trust value in [55] can be extracted from wireless signals between agents to verify the validity of an agent's identity (in a probabilistic sense) during a Sybil attack. This article builds most closely on the previous work in [20] and [55], which demonstrates how a stochastic trust variable $\alpha_{ij}(t)$ can be derived based on statistical analysis and local measurements. Using this trust variable, we can design weights that define a switching topology, allowing for arbitrarily small deviations in the final global error estimation of the platoon. Additionally, the controller benefits from this trust, enabling it to switch smoothly from Cooperative Adaptive Cruise Control (CACC) to Adaptive Cruise Control (ACC), maintaining safety during sudden changes.
The contributions of this paper are as follows:
\begin{enumerate}
\item Trust to opinion
\item Trust to controller
\item Trust to observer

\end{enumerate}


\section{Graph Representation}
Let $G = (V, E)$ be an undirected graph with $n$ nodes, where:
\begin{itemize}
    \item $V = \{1, 2, \dots, n\}$ is the set of nodes.
    \item $E \subseteq V \times V$ is the set of edges.
    \item The adjacency matrix $A$ of the graph $G$ is an $n \times n$ matrix where:
\end{itemize}
\begin{equation}
        A_{ij} = \begin{cases}
            1, & \text{if node } i \text{ connect} \text{ to node } j , \\
            0, & \text{otherwise}.
        \end{cases}
    \end{equation}
\subsection{\color{red}Virtual Graph Generation}

For each vehicle $j \in \{1, 2, \dots, n\}$, we generate a virtual graph $G^{(j)} = (V^{(j)}, E^{(j)})$ by adding an extra node $0$ and connecting it to node $j$ and its neighbors.

\begin{itemize}
    \item The set of nodes in the virtual graph $G^{(j)}$ is $V^{(j)} = \{0\} \cup V$.
    \item  The set of edges in the virtual graph $G^{(j)}$ is $E^{(j)} = E \cup \{(0, j)\} \cup \{(0, k) \mid k \in \text{neighbors of } j$
    \item The adjacency matrix $V^{(j)}$ of the virtual graph $G^{(j)}$ is an $(n+1) \times (n+1)$ matrix defined as:
    \begin{equation}
        V^{(j)} = \begin{bmatrix} 0 & \mathbf{v}_j^T \\ \mathbf{v}_j & A \end{bmatrix}
    \end{equation}
    where $\mathbf{v}_j$ is a column vector of size $n$ such that:
    \begin{equation}
        \mathbf{v}_j(i) = \begin{cases}
            1, & \text{if } i = j \text{ or } A_{ji} = 1, \\
            0, & \text{otherwise}.
        \end{cases}
    \end{equation}
    This means that node $0$ is connected to node $j$ and all its neighbors in the original graph.
\end{itemize}

% \subsection{Consensus Weights Calculation}
% For each virtual graph $G^{(j)}$, we calculate the consensus weight matrix $W^{(j)}$. The weight matrix $W^{(j)}$ is an $(n+1) \times (n+1)$ matrix where $W^{(j)}(i, l)$ represents the weight assigned to the edge between node $i$ and node $l$.

% The degree $d_i$ of node $i$ in $G^{(j)}$ is given by:
% \begin{equation}
%     d_i = \sum_{l=1}^{n+1} V_j(i, l)
% \end{equation}

% The weight matrix $W^{(j)}$ is then defined as:
% \begin{equation}
%     W_{i, l}^{(j)} = \begin{cases}
%         \frac{1}{d_i + 1}, & \text{if } V_j(i, l) = 1 \text{ or } i = l, \\
%         0, & \text{otherwise}.
%     \end{cases}
% \end{equation}

% This means that the weight $W_{i, l}^{(j)}$ is $\frac{1}{d_i + 1}$ if there is an edge between node $i$ and node $l$ or if $i = l$ (self-loop), and $0$ otherwise.

% \subsection{Summary}
% \begin{itemize}
%     \item \textbf{Virtual Graph $G_j$:} For each vehicle $j$, the virtual graph $G_j$ is constructed by adding node $0$ and connecting it to node $j$ and its neighbors.
%     \item \textbf{Weight Matrix $W_j$:} The weight matrix $W_j$ is calculated based on the degree of each node in the virtual graph $G_j$.
% \end{itemize}

% \subsection{General Form}
% For a general graph $G = (V, E)$ with adjacency matrix $A$, the virtual graph $G_j$ and its corresponding weight matrix $W_j$ can be expressed as:
% \begin{equation}
% \begin{aligned}
%     G_j = (V_j, E_j), \quad V_j = \{0\} \cup V, \\ E_j = E \cup \{(0, j)\} \cup \{(0, k) \mid k \in \text{neighbors of } j\}    
% \end{aligned}
% \end{equation}

% \begin{equation}
%     W_j(i, l) = \begin{cases}
%         \frac{1}{d_i + 1}, & \text{if } (i, l) \in E_j \text{ or } i = l, \\
%         0, & \text{otherwise}.
%     \end{cases}
% \end{equation}


\section{General description}\label{sec:genral}
\subsection{Local System}
The Minicar has Ackermann steering geometry, and its kinematics can be approximated by the bicycle model, with motion equations as follows:
\begin{align} \label{diff_1_LPV}
    \dot{X} &= v \cos(\theta) \\ \notag
    \dot{Y} &= v \sin(\theta)\\ \notag
    \dot{v} &= a \\ \notag
    \dot{a} &= \frac{1}{\tau}u - \frac{1}{\tau}a  \\ \notag
     \dot{\theta} &= L^{-1} v \tan(\delta)  \\
    % or \quad
    % \dot{\theta} &= \omega  
\end{align}
where
\begin{itemize}
    \item $X , Y$ denote the position,
    \item v is is the forward speed , a acceleration control 
    \item $\delta $ is the steering angle,
    \item L is the vehicle's wheel base. 
\end{itemize}

Each vehicle has a full measurement state obtained in real-time \textbf{position} and \textbf{velocity}, \textbf{heading angle} from a GPS, IMU sensor mounted on the vehicle. 

The radar sensor is located on the front of the
 host vehicle and measures the relative speed and distance to the rear bumper of the target
 vehicle. The relative distance $r$ is defined as

\begin{equation}\label{radar}
    r = \sqrt{(X_t - X_h)^2 + (Y_t - Y_h)^2} - L 
\end{equation}
$X$ and $Y$ positions of the host and preceding vehicle vehicle
 correspond with their center of gravity. Note that the expression in \eqref{radar} is a simplification
 where the difference in heading angle between the two vehicles is ignored. The effects of
 this simplification can be included in the measurement noise, which also accounts for the
 $r$ is defined as uncertainty of the measurement location on the target vehicle's bumper. The relative speed $\dot{r}$
 \begin{equation}\label{radar_speed}
     \dot{r} = (v_t - v_h) cos(\theta_r)
 \end{equation}
where $r$ is the azimuth angle measured by the radar. The relative speed is derived from a
 constant radius cornering assumption where the vehicles drive with the same curvature
 and have no relative lateral offset with respect to the desired trajectory.

%  \textcolor{red}{The dynamics of $i^{th}$ vehicle are a Position-Velocity third-order linear model given by the following set of equations:}

% \begin{equation}\label{eq_1}
%       \left\{ \begin{array}{l}
%         \dot{p}_i(t) =v_i(t) \\
%         \dot{v}_i(t) =a_i(t)  \\
%         \end{array}  \right.
% \end{equation}  
% where
% \begin{itemize}
% \item $p_i(t), v_i(t)$ and $a_i(t)$ denote the position, velocity and
% acceleration of the $i^{th}$  vehicle, respectively;
% \item $u_i(t)$ is the $i^{th}$ vehicle control input representing the desired
% driving/braking force or acceleration;
%  \item $\tau$ is the engine's constant time lag. 
 
%  \end{itemize}

%  All the equations can written :

% \begin{equation}\label{CACC}
% %\displaystyle
%        \begin{array}{l}
%         {x_i}(t+1)=A x_i + B  u_i + B_f f(x)_i  \\
%         y_{i}(t) = Cx_{i}\end{array} .
% \end{equation}  

% where 
% \begin{equation*}
%         A  =\begin{bmatrix}
%         0 & 1 & 0 \\
%         0 & 0 & 1 \\
%         0 & 0 & \frac{-1}{\tau}
%           \end{bmatrix},
%         B =\begin{bmatrix}
%         0 \\
%         0 \\
%         \frac{1}{\tau}
%           \end{bmatrix}, 
%         B_f =\begin{bmatrix}
%         0,0 \\
%         1,0 \\
%         0,1
%           \end{bmatrix},
%         C =\begin{bmatrix}
%         1 & 0 & 0 \\
%         0 & 1 & 0 \\
%         0 & 0 & 1
%           \end{bmatrix}, 
% \end{equation*}





\subsection{Model Discretization}
\textcolor{red}{Not finsh}

 To adopt a state-space notation , a cascading system description is considered, allowing us to transform the system
 dynamics into state-space form. The schematics of the cascading systems approach is
 shown in Figure \ref{fig:2sys},

\begin{figure*}[htp]
\hrule
    \centering
    \includegraphics[scale=0.5]{img/2system.png}
    \caption{Schematic of the cascading systems.}
    \label{fig:2sys}
\hrule
\end{figure*}

$$
\begin{gathered}
\mathbf{x}_1=\left[\begin{array}{ll}
  \theta_{\mathrm{h}} & \dot{\theta}_{\mathrm{h}}
\end{array}\right]^T \\
\mathbf{x}_2=\left[\begin{array}{lll}
 X_{\mathrm{h}} & Y_{\mathrm{h}} & v_{\mathrm{h}}
\end{array}\right]^T
\end{gathered}
$$

measurement vectors

$$
\begin{array}{c}
\mathbf{y}_1=\left[\begin{array}{ll}
 \theta_{\mathrm{h}} & \dot{\theta}_{\mathrm{h}}
\end{array}\right]^T \\
\mathbf{y}_2=\left[\begin{array}{lllll}
 X_{\mathrm{h}} & Y_{\mathrm{h}} & v_{\mathrm{h}} & \dot{\theta}_{\mathrm{h}} & r
\end{array} \right]^T \\
\end{array}
$$

and process noise vectors



are introduced. In state-space form, the system dynamics are described as

$$
\begin{aligned}
\dot{\mathbf{x}}_j=\mathbf{A}_j \mathbf{x}_j+\mathbf{B}_j \mathbf{u}_j+\mathbf{E}_j \mathbf{w}_j, \\
\mathbf{y}_j=\mathbf{C}_j \mathbf{x}_j+\mathbf{D}_j \mathbf{u}_j+\mathbf{J}_j \mathbf{v}_j,
\end{aligned}
$$

   
\subsection{Distributed/Collective system}


So we can write down the the virtual global system that represents all the k-neighbors of $i^th$ . 



\begin{equation}\label{CACC_dis_neibor}
        \mathbf{{X}}(t+1)=\mathbf{A} \mathbf{X} + \mathbf{B}({U + \mathcal{D}_{uk}}) 
\end{equation}

where 
\begin{equation*}
\mathbf{A} = I_N \otimes A_i, \quad \mathbf{B} = I_N \otimes B_i,
\end{equation*}




So the 

\begin{itemize}
\item $\mathbf{X} = [x_{i-k},..., x_{i} ,..., x_{i+k}]^T$ 
\item $U = [u_i , ... u_i ,.. u_i ]^T$ 


 \item $\mathcal{D}_{uk}$ is the uncertain in the controller input. 
 \end{itemize}


The Global System's state can be categorized into lead, follower, and last positions.


\section{ System Architecture and Design of controller}

\begin{figure*}[htp]
\hrule
    \centering
    \includegraphics[scale=0.5]{img/local_global_com.png}
    \caption{Schematic diagram of the estimation method.}\label{fig:information_boarcard}
\hrule
\end{figure*}

We will have two controllers: the inner controller (local control) and the outer controller (distributed control), located in the internal and external control loops of the vehicle, respectively. The inner controller is responsible for ensuring that the vehicle follows its reference trajectory while maintaining a safe distance from the preceding vehicle; this is commonly referred to as the Adaptive Cruise Control (ACC) controller. The outer controller, which we will refer to in this paper as the Distributed Observer Cooperative Adaptive Cruise Control (CACC), is designed to facilitate consensus within a platoon of vehicles. The CACC aims to minimize the distance between vehicles in a platoon so that the group can travel more quickly and smoothly. These 2 controllers will be combined to get the final controller to actually put in the control loop. 
\ref{fig:information_boarcard}



\subsection{Inner/Local controller}
\textcolor{red}{ Only use local estimation }

 The local system on longitudinal motion by using Intelligent Driver Model (IDM) was first proposed in $[..]$.
 The idea under-pinning IDM is that a vehicle's acceleration is a function of the vehicle's current velocity v, it's gap s
 to the preceding vehicle, and the approach rate $\Delta v$ to the preceding vehicle. It is formalized as :
 \begin{equation}
    a_{\mathrm{IDM}}=\alpha\left[1-\left(\frac{v}{v_0}\right)^\delta-\left(\frac{s^{\star}(v, \Delta v)}{s}\right)^2\right]
 \end{equation}

where $s^{\star}$ is a function determining the desired minimum gap to the preceding vehicle. We compute this value as

\begin{equation}
s^{\star}(v, \Delta v)=s_0+T v+\frac{v \Delta v}{2 \sqrt{\alpha \beta}}.    
\end{equation}

where $s_0$ corresponds to a jam distance. 
\textcolor{red}{ Only use global estimation}
\subsection{Cooperative controller}
Nevertheless, the classical IDM that only considers the impacts of the nearest preceding vehicle is insufficient to describe the CAV
behaviors, where each CAV exchanges information with multiple surrounding vehicles. Hence, a revised IDM incorporating cooperation through wireless communications, labeled cooperative IDM (CIDM) \cite{C_IDM}, is formulated by the following equations:
\begin{equation}\label{CIDM}
a_{CIDM}=\alpha \left[1-\left(\frac{v_n}{v_0}\right)^4-\left(\frac{S^{*}\left(v_n, \Delta \tilde{v}_{n-j+1}\right)}{\sum_{j=1}^N \alpha_{n j} w_{n j} \tilde{{s}}_{n-j+1}}\right)^2\right] 
\end{equation}
with
\begin{equation*}
\begin{aligned}
    S^{*}\left(v_n(t), \Delta \widetilde{v}_{n-j+1}(t)\right)=s_0+T v_n(t)+ \\ \frac{v_n(t) \cdot \sum_{j=1}^N \beta_{n j} w_{n j} \Delta \tilde{v}_{n-j+1}(t)}{2 \sqrt{a b}} .
\end{aligned}
\end{equation*}

In the steady-state, all vehicles travel at the same velocity $\nu^{\text {eq }}$ with the same clearance gap $s^{\text {eq }}$. Eq. \eqref{CIDM} yields the uniform clearance gap as follows:

\begin{equation}\label{gap_CIDM}    
s^{\mathrm{eq}}(v)=\left(s_0+T v^{\mathrm{eq}}\right)\left(1-\left(\frac{v^{\mathrm{eq}}}{v_0}\right)^4\right)^{-1 / 2}
\end{equation}
Eq. \eqref{gap_CIDM} shows that each following vehicle's gap at equilibrium is determined by the equilibrium velocity. However, once an
attack occurs, the equilibrium clearance gap and velocity will be perturbed.

\subsection{Final Controller}
The target to the vehicle is 
\begin{equation}
    a_{target} = (1-\gamma). a_{IDM} + \gamma.a_{CIDM} 
\end{equation}
where, the parameter $\gamma$ is used for tuning the control strategy, can be determined by the level of trust we have in the connection.. It allows us to combine two types of control for optimal performance. When $\gamma = 0$, we rely entirely on cooperative control, which is beneficial when we have a strong connection. Conversely, when $\gamma = 1$, we use the Intelligent Driver Model (IDM) solely based on local sensor data to drive the controller. 

Apply the First-Order Filter
\begin{equation}
    u_{t} = u_{t-1} + \tau_{f}(a_{target} - u_{t-1})
\end{equation}

The vehicle's control input at time $t$ is computed through the integrator $v_t= u_{t}$. $\Delta t + v_{t-1}$.

%% ------------------------------------

\section{Distributed Driving State Estimation Scheme}
Similarly to controller, the estimation step also has 2 distinct designs. 
\subsection{Local Observer Design}
\begin{equation}\label{observer_dis_plug}
\bar{x}_i(k+1)=A \bar{x}_i(k)+B u_i(k)+F_i(y_i(k)-C \bar{x}_i(k))
\end{equation}
The estimation error is : 
\begin{equation}
e_0^{(j)}(k+1)= A - F_j C
\end{equation}
The main goal is to find the appropriate values for the observer gains $F_j$ by solving certain LMIs [..] or Kalma filter. This is done in order to ensure that the estimation error $e_0^{(j)}$ converges to zero asymptotic. 

\begin{remark}
In a local observer, we have the flexibility to develop a better model for improved estimation results, such as using model [...]. The idea is that a local observer will consistently provide accurate estimates, while remaining decoupled from the distributed observer.
\end{remark}

\subsection{Distributed Observer Design}
In this section, we focus on a collective layer system in which each vehicle can estimate the platoon's stats, so we call this distributed observer. We develop a distributed driving state observer for a vehicle platoon and establish conditions that guarantee
the stability of its estimation error dynamics.

We begin by presenting an overview of our design. Based
on the observability analysis of the platooning system, each
vehicle $i \in V$ is able to estimate its own state $x_i(k)$ using
local observer and its local measurement $y_i(k)$. The estimate is denoted as $\bar{x}_i(k)$ and is updated by vehicle $i$ 
locally. The estimate $\bar{x}_i(k)$ is then treated as an external signal
to be tracked by all vehicles in the platoon using a leader-follower consensus technique. This allows each vehicle $i$ to
construct an estimate of vehicle $j's$ state by tracking the signal
$\bar{x}_j(k)$ generated by vehicle $j$ using its own local observer for $\forall j \in V$. We represent the estimate of vehicle $j$
as $\hat{x}_i^{(j)}(k)$ by vehicle $i \in V$. Finally, each vehicle $i$ develops
an estimate of the true global driving state $x(k)$, denoted as
$\hat{x}_i(k) = [(\hat{x}_i^{(1)}(k))^\top, \ldots, (\hat{x}_i^{(m)}(k))^\top]^\top$.
This distributed state estimation \cite{Plug_play} procedure can be summarized as follows :
% \begin{equation}
% \begin{aligned}
% &\bar{x}_i(k+1)=\left\{\begin{array}{l}
% A \bar{x}_1(k)+B u_1(k)+F_1\left(y_1(k)-C_{1,1} \bar{x}_1(k)\right), i=1 \\
% A \bar{x}_i(k)+B u_i(k)+F_i\left(y_i(k)-\left(C_{i, i} \bar{x}_i(k)\right.\right. \\
% \left.\left.+C_{i, i-1} \hat{x}_i^{(i-1)}(k)\right)\right), i \in \mathrm{~V} /\{1\}
% \end{array}\right.\\
% &\begin{aligned}
% \hat{x}_i^{(j)}(k+1)= & A\left(\hat{x}_i^{(j)}(k)+\sum_{l \in \mathrm{~N}_i} w_{i l}^{(j)}\left(\hat{x}_l^{(j)}(k)-\hat{x}_i^{(j)}(k)\right)\right. \\
% & \left.+w_{i 0}^{(j)}\left(\bar{x}_j(k)-\hat{x}_i^{(j)}(k)\right)\right)+B u_j(k), i, j \in \mathrm{~V}
% \end{aligned}
% \end{aligned}
% \end{equation}

\begin{equation}\label{dis_obs}
\begin{aligned}
\hat{x}_i^{(j)}(k+1)= & A\left(\hat{x}_i^{(j)}(k)+\sum_{l \in \mathrm{~N}_i} w_{i l}^{(j)}\left(\hat{x}_l^{(j)}(k)-\hat{x}_i^{(j)}(k)\right)\right. \\
& \left.+w_{i 0}^{(j)}\left(\bar{x}_j(k)-\hat{x}_i^{(j)}(k)\right)\right)+B u_j(k), i, j \in V
\end{aligned}
\end{equation}

where $w_{il}^{(j)}$ and
$w_{i0}^{(j)}$ are nonnegative consensus weights related to the communication
network $G$. These parameters will be determined in Section \ref{Weigth_Design}. 

\remark{It is difficult to get, or send the control input because it will be private. So we will consider that all the controllers in vehicle $j^{th}$ that is the neighbor of $i^{th}$ vehicle will be  $u_j = u_i + d_{u_j}$ , $d_{u_j}$ is a bound have a bound value. so the controller in the virtual global system  $U = [u_{1} , u_2 ,.. u_{j} ]^T$ }

\textbf{Explain the concept of an observer:} The system aims to get all states of the platoon, so the second term focuses on improving the global estimate in each vehicle by compensating with communication data. However, since the initial global state of the platoon in each vehicle is incorrect and we do not have precise information about the controllers, nor good models of other vehicles for accurate predictions, we address this uncertainty with the third term and use local estimation as a ref to achieve. The first term is then combined with the other terms to predict the overall global state of the system in the next step.

% \textbf{Explain the concept of an observer:} The system aims to reach a consensus value, so the second term focuses on improving the global estimate in each vehicle by compensating with communication data. However, since the initial global state of the platoon in each vehicle is incorrect and we do not have precise information about the controllers of other vehicles for accurate predictions, we address this uncertainty with the third term, which attempts to achieve a local estimation of these others. The first term is then combined with the other terms to predict the overall global state of the system in the next step.


\remark{Note that the information transmitted
by each vehicle $i \in V$ through the communication network $G$ is
the local estimate and the global estimate state of the platoon from each vehicle, i.e., $\hat{x}_i(k)$ and $\bar{x}_i(k)$.}

% \remark{k-neighbor of $i^th$ will send the local estimated value $\hat{x}_k $ to the $i^th$ vehicle and, also send about global estimated of all vehicles in the platoon.
% }

\section{Trust Score Framework (Modeling Trust)}



To prevent collisions induced by data injection attacks as demonstrated in \cite{C_IDM}, only accurate data can be passed to the cruise controller. Since a host vehicle cannot judge the value of a single data point received in a beacon, it has to be compared to the sender's actual behavior over a longer period of time. This allows the host vehicle to gain trust in the sender's benignity or detect misbehavior. Depending on the trust in the predecessor within the platoon, the host vehicle can regulate the safety gap to that vehicle. If the trust falls below a certain threshold, it might ultimately decide not to be in a platoon with that vehicle at all. This section introduces the details of TriP, our trust-based misbehavior detection system for platoons.


We are given the following trust-related information:
\begin{itemize}
    \item \textbf{Self Trust}: Each vehicle \( i \) has a trust score of 1 for its own local estimate \( \bar{x}_i(k) \).
    \item \textbf{Neighbor Trust}: Vehicle \( i \) evaluates trust scores for the local and global estimates of its nearby neighbors \( l \in \mathrm{N}_i \).
    \item \textbf{Trust Sharing}: Neighbors send their trust scores to vehicle \( i \).
    \item \textbf{Trust Aggregation}: Through aggregation, vehicle \( i \) computes trust scores for all vehicles in the platoon, not just its immediate neighbors.
\end{itemize}

Let's define:
\begin{itemize}
    \item \( t_{i}^{(l)} \): Vehicle \( i \)'s trust in vehicle \( l \) (for \( l \in \mathrm{N}_i \), this is directly evaluated; for others, it's aggregated).
    \item \( t_{l}^{(j)} \): Vehicle \( l \)'s trust in its estimate of vehicle \( j \)'s state, \( \hat{x}_l^{(j)}(k) \), which is sent to vehicle \( i \) if \( l \in \mathrm{N}_i \).
    \item \( t_{i}^{(i)} = 1 \): Vehicle \( i \) fully trusts its own local estimate.
\end{itemize}

\begin{figure*}[h!]
\hrule
    \centering
    \includegraphics[scale=0.5]{img/trust_framwork.png}
    \caption{Trust framework.}
    \label{fig:trust_framwork}
\hrule
\end{figure*}

\assumption{In the virtual global system, there naturally aren't faults; the faults are caused by uncertainties or cyber-attacks in the communication channels.}

Since the observer model uses two types of data the vehicle receives, we need to quantify the trustworthiness in these data. 

\subsection{Interactions of the local data}

A host vehicle with automated driving functions, such as Cooperative Adaptive Cruise Control (CACC), must ultimately trust its own sensor data. When assessing other platoon member's trustworthiness, the host vehicle uses its sensor data as ground truth to compare the others actual behavior to what they announced in their beacons. This comparison is feasible only for vehicles within the perception range of the host vehicle's local sensors and is most effective for the direct predecessor and successor. Inspire of \cite{trust_simple, trust_complex} each criterion evaluation results in a value between 0 (worst) and 1 (best). 

\paragraph{Velocity} The host vehicle calculates the reference velocity $v_\text{ref}$ based on the leader's last reported velocity $v_L$, adjusted by the reported acceleration $a_L$ and the elapsed time $b_L$ since the last leader beacon:
\begin{equation}
    v_\text{ref} = v_L + b_L \cdot a_L.
\end{equation}

\begin{itemize}
    \item Target Mismatch with the Leader
        \begin{equation}
        v_{y,\text{ref}} = 
        \begin{cases} 
            \max\left(1 - \left|\frac{v_y - v_\text{ref}}{v_\text{ref}}\right|, 0\right), & \text{if } v_\text{ref} > 0, \\
            \max\left(1 - |v_y|, 0\right), & \text{otherwise}.
        \end{cases}
    \end{equation}
    \item Target Mismatch with the Host
        \begin{equation}
            v_{y,\text{host}} = 
            \begin{cases} 
                \max\left(1 - \left|\frac{v_y - v_\text{host}}{v_\text{host}}\right|, 0\right), & \text{if } v_\text{host} > 0, \\
                \max\left(1 - |v_y|, 0\right), & \text{otherwise}.
            \end{cases}
        \end{equation}
    % \item Host Mismatch with the Leader
    %     \begin{equation}
    %         v_{host,\text{leader}} = 
    %         \begin{cases} 
    %             \max\left(1 - \left|\frac{v_{host} - v_\text{ref}}{v_\text{ref}}\right|, 0\right), & \text{if } v_\text{ref} > 0, \\
    %             \max\left(1 - |v_y|, 0\right), & \text{otherwise}.
    %         \end{cases}
    %      \end{equation}
\end{itemize}
Fusion of these Mismatches using weighted average :
\begin{equation}
    v_{y,final} = \sigma. v_{y,leader} + (1-\sigma). v_{y,host}     
\end{equation}

The alpha will change depending on the position of the target vehicle. If the target is following the host, then alpha will ideally be around 0-0.3. If the target is ahead of the host, alpha will be around 0.7-1.

\paragraph{Distance} The host vehicle calculates the error between the measured distance $d$ from its sensors and the reported distance $d_{b,y}$ from the preceding vehicle:
\begin{equation}
    d_y = \max\left(1 - \left|\frac{d_{b,y} - d}{d}\right|, 0\right).
\end{equation}

\paragraph{Acceleration} The relative velocity $v_\text{rel}$ is derived as the rate of change in distance $d$ over the sample interval $t_s$. The acceleration criterion is:
\begin{equation}
    a_y = \max\left(1 - \left|\frac{v_\text{rel}}{t_s} \cdot a_\text{diff}\right|, 0\right),
\end{equation}
where $a_\text{diff}$ is the difference between the received and measured host acceleration. The term $\frac{v_\text{rel}}{t_s}$ scales the effect of acceleration mismatch based on how fast the vehicles are moving relative to each other.
Acceleration differences matter more if two vehicles are close together and moving at different speeds.
If they are farther apart or moving at the same speed, acceleration differences are less critical.
This prevents small acceleration mismatches from being overly penalized when relative motion is low.

% Why is the Equation Designed This Way?
% A. Measuring Acceleration Consistency
% The trust formula measures how much the reported acceleration deviates from the actual acceleration.
% If the difference is large, it means the vehicle is not reporting correct acceleration values, which could be due to:
% Malicious behavior (e.g., injecting false data).
% The larger the deviation, the more trust should decrease.
% B. Scaling the Trust Impact Based on Relative Motion
% The term scales the effect of acceleration mismatch based on how fast the vehicles are moving relative to each other.
% Acceleration differences matter more if two vehicles are close together and moving at different speeds.
% If they are farther apart or moving at the same speed, acceleration differences are less critical.
% This prevents small acceleration mismatches from being overly penalized when relative motion is low.

%------------Jerkiness
% \paragraph{Jerkiness:} The absolute jerk $j_{y,\text{abs}}$ is evaluated, with penalties applied only if it exceeds a threshold $j_\text{thres}$:
% \begin{equation}
%     j_y = \min\left(\frac{j_\text{thres}}{j_{y,\text{abs}}}, 1\right).
% \end{equation}

\paragraph{Beacon Delivery Timeout} If a member vehicle $y$ fails to send a beacon within the expected timeout, the timeout factor $t_y$ is set to 0, otherwise it is set to 1.

\paragraph{Trust Samples}

All criteria are combined into a single trust sample $t_y$:
\begin{equation}
    t_y = t_{o,y} \cdot v^{w_v}_{y,\text{ref}} \cdot d^{w_d}_y \cdot a^{w_a}_y 
\end{equation}
% \begin{equation}
%     t_y = t_{o,y} \cdot v^{w_v}_{y,\text{ref}} \cdot d^{w_d}_y \cdot a^{w_a}_y \cdot j^{w_j}_y,
% \end{equation}
where $w_v, w_d, w_a$ are the respective weights for the criteria. So the maximal value of trust will be : 
$\max{t_y} \leq 1$


\textcolor{red}{Big Question: We only use local to evaluate trust , how if the global being attack ? }

\textcolor{blue}{Answer: we need to check again the global ?, maybe not like local, but something like mean and variance? After that, by magic somehow can check the global, we could use that to change the weight, so that will related to parameter $\kappa$ in \eqref{kappa} }


\subsection{Interactions of the global data}


\paragraph{Consistency with Host's Global Estimate}


To assess whether neighbor \( l \)'s global estimate is trustworthy, the host can compare it directly to its own global estimate. Here's how this can be done step-by-step:
\assumption{We assume that the change in state a between any two consecutive time instances $k$ and $k + 1$ to be approximately in a small bound. This is generally true of all the vehicle in practice, to ensure the movement of the vehicles to be smooth} 
\begin{enumerate}
    \item We use the \textbf{Mahalanobis distance}, which normalizes discrepancies based on covariance:
\[
d_{i,l}^{(j)}(k) = \left( \mathbf{\hat{x}}_i^{(j)}(k) - \mathbf{\hat{x}}_l^{(j)}(k) \right)^T \mathbf{\Sigma}^{-1} \left( \mathbf{\hat{x}}_i^{(j)}(k) - \mathbf{\hat{x}}_l^{(j)}(k) \right)
\]
where:
\begin{itemize}
    \item \( \mathbf{\hat{x}}_i^{(j)}(k) \) is the host's estimate of vehicle \( j \) (e.g., position, velocity, acceleration).
    \item \( \mathbf{\hat{x}}_l^{(j)}(k) \) is the neighbor's estimate of vehicle \( j \).
    \item \( \mathbf{\Sigma} \) is the \textbf{covariance matrix}, which captures variance and correlations between state variables.
\end{itemize}

The \textbf{the covariance matrix } is diagonal:
\[
\mathbf{\Sigma} = \text{diag}(\sigma^2_p, \sigma^2_v, \sigma^2_a)
\]
where:
\begin{itemize}
    \item \( \sigma^2_p \) is the variance of position discrepancies.
    \item \( \sigma^2_v \) is the variance of velocity discrepancies.
    \item \( \sigma^2_a \) is the variance of acceleration discrepancies.
\end{itemize}





    \item \textbf{Sum the Discrepancies Across the Platoon}  
    Add up these differences to get a total discrepancy score:
    \[
    D_{i,l}(k) = \sum_{j=1}^n d_{i,l}^{(j)}(k)
    \]
    A small \( D_{i,l}(k) \) means \( l \)'s global estimate is close to the host's, suggesting higher trustworthiness.

    \item \textbf{Convert to a Trust Factor}  
    Turn this discrepancy into a trust score using an exponential function:
    \[
    \gamma_{i,l}^{\text{host}}(k) = \exp\left( -{D_{i,l}(k)} \right)
    \]
    \begin{itemize}
        \item The result, \( \gamma_{i,l}^{\text{host}}(k) \), ranges from 0 to 1. A value near 1 means high trust (small discrepancy), while a value near 0 means low trust (large discrepancy).
    \end{itemize}
\end{enumerate}

\begin{remark}
    
\begin{itemize}
    \item If velocity has high variance, \( \sigma^2_v \) will be large, making velocity discrepancies contribute less to the final distance.
    \item If position has low variance, \( \sigma^2_p \) will be small, making position discrepancies contribute more.
    \item If acceleration and velocity are correlated, Mahalanobis distance accounts for this, \( \mathbf{\Sigma} \) will contains off-diagonal terms capturing dependencies.
\end{itemize}
\end{remark}



\paragraph{Consistency with Local Sensor Measurements}\label{local_sensor}

\textbf{Objective}: Ensure that global estimates align with trusted local sensor data for nearby vehicles.

\textbf{Approach}: Vehicle \( i \) uses its sensors to measure the relative state of nearby vehicles (e.g., distance or relative velocity to a predecessor) and checks if the global estimates from neighbor \( l \) are consistent with these measurements.

\textbf{Mathematical Implementation}:

\begin{itemize}
    \item \textbf{Local Measurement}: Let \( y_{i,j}(k) \) be vehicle \( i \)'s sensor measurement of vehicle \( j \)'s state relative to itself (e.g., \( y_{i,j}(k) = x_j(k) - x_i(k) \) for relative position). This is available for vehicles \( j \in M_i \), where \( M_i \) is the set of vehicles within \( i \)'s sensor range.

    \item \textbf{Expected Consistency}: For neighbor \( l \)'s global estimate to be plausible, it should satisfy:
    \[
    \hat{x}_l^{(j)}(k) - \hat{x}_l^{(i)}(k) \approx y_{i,j}(k)
    \]

    \item \textbf{Consistency Error}: Calculate the error between the global estimate and the local measurement:
    \begin{align}
    e_{i,l}^{(j)}(k) = \left( \hat{x}_l^{(j)}(k) - \hat{x}_l^{(i)}(k) \right) - y_{i,j}(k) \\
    E_{i,l}^{(j)}(k) = (e_{i,l}^{(j)}(k))^{-1}. \mathbf{\Sigma}^{-1}. e_{i,l}^{(j)}(k)
    \end{align}

    \item \textbf{Total Consistency Error}: Sum the errors over all measurable vehicles:
    \[
    E_{i,l}(k) = \sum_{j \in M_i} E_{i,l}^{(j)}(k)
    \]
    A higher \( E_{i,l}(k) \) suggests that neighbor \( l \)'s global estimate contradicts local sensor data.

    \item \textbf{Trust Factor}: Convert the error into a trust factor:
    \[
    \gamma_{i,l}^{\text{local}}(k) = \exp\left( -\frac{E_{i,l}(k)}{\tau^2} \right)
    \]
    \begin{itemize}
        \item Range: \( \gamma_{i,l}^{\text{local}}(k) \in (0, 1] \), with lower values indicating poor consistency.
    \end{itemize}
\end{itemize}

This trust factor ensures that global estimates align with reliable local observations.

\subsection{Combine with Other Trust Checks}  
\textcolor{red}{ !!! Not sure , why take half global , half local , something else better ?}
    The host likely already uses other ways to evaluate trust, like:
    \begin{itemize}
        \item Checking if \( l \)'s estimate matches the host's local sensor data (a \textbf{local consistency} factor, \( \gamma_{i,l}^{\text{local}}(k) \)).
        % \item Comparing \( l \)'s estimate with what other neighbors say (a \textbf{cross-validation} factor, \( \gamma_{i,l}^{\text{cross}}(k) \)).
    \end{itemize}
    We can combine these into an overall trust score:
    \begin{equation}\label{trs_sample_genral}
    t_{i,l}^{\text{ext}}(k) = \frac{1}{2}(t_y +  \gamma_{i,l}^{\text{host}}(k) \cdot \gamma_{i,l}^{\text{local}}(k) 
    )    
    \end{equation}

    or
\begin{equation}\label{trs_sample_genral}
    t_{i,l}^{\text{ext}}(k) = t_y \cdot  \gamma_{i,l}^{\text{host}}(k) \cdot \gamma_{i,l}^{\text{local}}(k)     
\end{equation}
    
    This extended trust score reflects how well \( l \)'s global estimate aligns with the host's view and its own sensors

\begin{remark}
    While we will proceed with trust, we must not forget the importance of prior calculations. This means we can flag anomalies when something appears to be incorrect. For instance, in the calculation of global trust, we consider the local gamma and cross gamma. If either of these values is abnormally low, we will raise a flag to the use after.
\end{remark}




\subsection{Building general trust for nearby vehicle}
Since the resulting trust samples \eqref{trs_sample_genral} for interactions may change over time and a trust is build over repeated interactions, these values need to be aggregated in a way that allows historic evolution, judgment on reliability and
certainty. Thus, the aggregation can be represented in a probability distribution function (PDF) by using the framework
of Bayesian statistics.


Bayesian trust systems from literature are often based on the Beta Reputation System [21] which utilizes the beta PDF [15]. The Beta Reputation System is only able to express ratings in two outcomes, such as trusted and untrusted. The Dirichlet Reputation System [22] however
enables more fine granular ratings compared to the Beta Reputation System, as is is based on the Dirichlet PDF which is the generalization of the Beta PDF, supporting an arbitrary number of outcomes.


To describe the trustworthiness, k = 5 mutually exclusive, equidistant and ascending trust rating levels [01] are defined as \textit{untrustworthy, bad, acceptable, good } and \textit{excellent}.

When vehicle x calculates a $t_y \in [0, 1]$ for vehicle y, it is mapped to a $k$-dimensional vector $r_y^x$ where every element of the vector is 0 except the one for the respective trust rating level (which is set to 1). The vector $r_y^x$ can be interpreted as probability distribution with a single outcome, determined
by the trust sample. 

Multiple outcome vectors $r_y^x$ from x for y are then aggregated in the rating vector:
\begin{equation}
    R_y = (R_y(1), . . . , R_y(k))
\end{equation}


Trust samples are aggregated over time using a Bayesian approach. The aggregated rating vector $R_y$ accumulates past outcomes as follows:
\begin{equation}
    R_{y,(t+1)} = (1 - \lambda_y) \cdot R_{y,t} + r_{xy},
\end{equation}
where $\lambda_y$ is an aging factor depending on the trust score $T_i$:
\begin{equation}
    \lambda_y = T_i \cdot w_t.
\end{equation}
where $w_t = 1$ mean quick change the trust , $=0$ mean slow change the trust  
\textcolor{red}{where $w_t$ is higher like $(0.7 - 0.8)$ when vehicle perform only longitudinal , mean if a misbehavior detected, will decease fast the trust. But in case change lane , $w_t = ( 0.4 - 0.6)$ some "positive false" detected misbehavior will lead change the trust less}

\textcolor{blue}{Note: positive false" detected misbehavior, meaning that: a vehicle can think the neighbor is attacker, but actually no , since when changelane sometime not too smooth}


The probability distribution vector $S_y$ is computed using an a priori distribution vector $a$:
\begin{equation}
    S_y(i) = \frac{R_y(i) + C \cdot a(i)}{C + \sum_{j=1}^k R_y(j)}.
\end{equation}

The trust score $\alpha$ is then calculated as:
\begin{equation}
    T_y = \sum_{i=1}^k \frac{i-1}{k-1} \cdot S_y(i).
\end{equation}

\textcolor{red}{Question : How to get the trust value of the other vehicle that we don't have direct measurement to compare , to build the Trust ?} 
\subsection{Building General Trust  for all vehicle in platoon}

We summarize trust in 3 case : 
\paragraph{Case 1 : Self-Trust }
\begin{equation}
    \mathcal{O}_i(i) = T_i(i) = 1
\end{equation}

\paragraph{Case 2 : Direct Nearby Vehicle}
\begin{itemize}
    \item For vehicles that are nearby neighbors (denoted $j \in N_i$, where $N_i$ is the set of neighbors of vehicle $i$), vehicle $i$ can evaluate their trust directly using local measurements (e.g., sensor data, communication reliability, etc.).
    \item This direct trust score is denoted $t_i^{(j)}$.
    \item \textbf{Key Point}: Since this trust is based on vehicle $i$'s own measurements, it's reliable and doesn't require aggregation. Thus, the trust score for these neighbors is simply:
    \begin{equation}
        \mathcal{O}_i(j) = T_i(j) = t_i^{(j)}
    \end{equation}
    \item No aggregation is needed here; the vehicle $i$ maintains the trust that it evaluates directly.
\end{itemize}

% This reflects full trust in its own estimate, and no aggregation or external input is involved.

\paragraph{Case 3 : Neighbors Without Direct Measurement and Non-Neighbors}

\item For the last case:
    \begin{itemize}
        \item \textbf{Neighbors without direct measurement}: If there's a neighbor $j \in N_i$ for which vehicle $i$ cannot compute a direct trust score $t_i^{(j)}$ (e.g., due to missing sensor data).
        \item \textbf{Non-neighbors}: Vehicles $j \notin N_i$ that are not in the neighbor set of vehicle $i$, so no direct measurement is possible.
    \end{itemize}
In both cases, vehicle $i$ cannot rely on its own evaluation and instead aggregates trust scores shared by its neighbors.
    

In a vehicular trust model, a vehicle \( i \) can directly evaluate the trust of its neighboring vehicles (front and back). However, for other vehicles \( j \), vehicle \( i \) does not have a direct trust score but can use trust scores shared by other vehicles \( k \) that have direct observations of \( j \). This concept is known as \textbf{opinion-based trust aggregation}.

We define the \textbf{opinion score} as a weighted sum of trust values received from other vehicles:

\begin{equation}
\mathcal{O}_i(j) = \sum_{k \in K_j} w_{ik} \cdot T_k(j)    
\end{equation}

where:
\begin{itemize}
    \item \( T_i(j) \) is the estimated trust score of vehicle \( j \) from vehicle \( i \)'s perspective.
    \item \( T_k(j) \) is the trust score of vehicle \( j \) provided by another vehicle \( k \) (which has direct observation of \( j \)).
    \item \( K_j \) is the set of vehicles that provide trust values for vehicle \( j \). $K_j$ should be equal to $N_i$ (neighbor of $i$)
    \item \( w_{ik} \) is the \textbf{weight} assigned to each vehicle \( k \), based on its distance from \( j \).
\end{itemize}

\subsection{Choice of Weighting Function}
We can have 2 main type to have the weight for opinion score :  
\paragraph{Weight based distance for $\mathcal{O}_{dis}$}
To determine \( w_{ik} \), we use an \textbf{inverse distance-based weighting function}, meaning that trust scores from closer vehicles are given higher importance:

\[
w_{ik} = \frac{1}{1 + |k - j|}
\]

where:
\begin{itemize}
    \item \( |k - j| \) is the absolute difference in vehicle indices in the platoon.
    \item The denominator ensures that closer vehicles receive a higher weight.
    \item These weights are then \textbf{normalized} so that they sum to 1:

    \[
    w_{ik} = \frac{w_{ik}}{\sum_{m \in K_j} w_{im}}
    \]
\end{itemize}

This ensures that the final opinion score is a properly weighted trust estimate.\ref{fig:opinon_trust}
\paragraph{Weight based Trust propagation for $\mathcal{O}_{popa}$}

\begin{equation}
    % w_{il} = \frac{t_i^{(l)}}{\sum_{m \in N_i} t_i^{(m)}}
    w_{ik} = \frac{t_{i}^k}{ |K_j|}
\end{equation}
where:
\begin{itemize}
    \item For any vehicle $j$ (including neighbors and non-neighbors), each neighbor $k \in N_i$ provides $t_k^{(j)}$, its trust in its estimate of $j$. 
    \item Vehicle $i$ trusts $k$ with score $t_i^{(k)}$, so the combined trust $i$ has in $k$'s estimate of $j$ is $t_i^{(j)} =  t_i^{(k)} \cdot t_k^{(j)}$. 
\end{itemize}

Exemple : 	In vehicle 3 , we have $\alpha(0) = [1 1 1 1]$ , and $\alpha(k) = [....]$

with $\alpha(k) = [. , 0.9, 1, 0.9]$ can caluated
If we want to calculate $\alpha_1$ that is for vehicle 1. use the received alpha to do that , like $(0.8 \cdot 0.9 )$  or $(0.5 \cdot 0.9 )$ , sum it up and divide by number data we use $(here = 2)$

But in case 5 vehicle , same procedure !!! 

Need to get a formula mathematic



    
\begin{figure}[htp]
    \centering
    \includegraphics[scale=0.45]{img/opinon_trsut.png}
    \caption{New trust (Opinion). Trust Propagation Mechanism}
    \label{fig:opinon_trust}
\end{figure}


After take the mean opinion of 2 type of weight (distance and popagation):
\begin{equation}
\mathcal{O}_i(j) = \frac{\mathcal{O}_{dis} + \mathcal{O}_{popa}}{2}    
\end{equation}



\section{Weigth Design for observer}\label{Weigth_Design}


\subsection{Choice of Weights}
Since we can find which vehicle is legitimate and which is malicious. We can design a new Weight matrix to prioritize the value of legitimate and forget the malicious value. To do so 
First, we redefine a time-dependent trusted neighborhood for each agent $i$ as

\begin{equation*}
\mathcal{N}_i(t)=\left\{j \in \mathcal{N}_i: \mathcal{O} > 0.5 \right\}    
\end{equation*}

That mean the connection is described in good or excellent.

This is the set of neighbors that legitimate agent $i$ classifies as its legitimate neighbors at time $t$. Denote

\begin{equation}\label{new_setneibor}
n_{w_i}^{G(j)}(t)=\max \left\{\kappa,\left|\mathcal{N}_i^{G(j)}(t)\right|+1\right\} \geq 1 \quad \text { for all } i \in \mathcal{L}
\end{equation}


where $\kappa>0$ can be thought of as a parameter limiting the maximum influence that the neighbors of agent $i$ are allowed to have on agent $i$ 's update values (i.e., a value of $\kappa$ approaching infinity would preclude neighbors from influencing agent $i$ 's updated values). Using these quantities, we define the weight matrix $W(t)$ by choosing its entries $w_{i j}(t)$ as follows: for graph $G^{(j)}$ with every $i\in L$, $l \in N_i$, we could using 2 design pattern for the weight following :


\textbf{ The  weight matrix that priority the local estimation term is defined as }:
\begin{equation}\label{weight_0}
w_{i l}^{(j)}(t)= \begin{cases}
\frac{1}{n_{w_i}^{G(j)}(t)}, & \text { if } l \in \mathcal{N}_i(t) \cup\{0\} \\ 
1-\sum_{m \in \mathcal{N}_i} w_{i m}(t), & \text { if } l=i \\
0, & otherwise \\ 
\end{cases}
\end{equation}

\textcolor{blue}{We can have another weigh update by using of Weighted Degree and Normalization like following .}


To ensure \( W(t) \) is sub-stochastic, we need to normalize the weights such that the sum of each row does not exceed 1.

\textbf{Step 1: Compute Unnormalized Weights}
For each node \( i \), compute the unnormalized weights based on the trust values:
\[
\tilde{W}_{il}(t) = \begin{cases}
f(\tau_{il}(t)) & \text{if } (i, l) \in E_j(t) \text{ or } i = l, \\
0 & \text{otherwise}.
\end{cases}
\]
Here, \( f(\tau_{il}(t)) \) is the trust-to-weight mapping function (e.g., linear mapping \( f(\tau) = \tau \)).

\textbf{Step 2: Normalize the Weights}
Normalize the weights so that each row sums to at most 1:
\[
W_{il}(t) = \frac{\tilde{W}_{il}(t)}{\sum_{k=1}^{n+1} \tilde{W}_{ik}(t)}.
\]
This ensures:
\[
\sum_{l=1}^{n+1} W_{il}(t) = 1 \quad \forall i \in \{1, 2, \dots, n+1\}.
\]
If the sum of unnormalized weights is already less than 1, the normalization step will preserve this property.

% ----------------------------
% \textcolor{red}{Think about using Opinion here in to the weight , if we want a separate weight for $w_{ii}$}
% \textcolor{blue}{ The design gain need way more complex than now}
% \textbf{The  weight matrix that is distributed equally is defined as }:


% \begin{equation}\label{weight_equal}
% w_{i l}^{(j)}(t)= \begin{cases}
% \frac{1}{n_{w_i}^{G(j)}(t)}, & \text { if }, l \in \mathcal{N}_i^{G(j)}(t) \\ 
% 0, & otherwise
% \end{cases}
% \end{equation}

% \textbf{New  weight matrix }:

% \textcolor{blue}{ For nearby vehicle , we will use flag for the weight , and for no near we use the cut set neigborh connected vehicle in \eqref{new_setneibor}  }

% \begin{equation}\label{weight_equal}
% \begin{cases}
% \text{with } l \in \mathcal{N}_i^{G(j)}(t) ,w_{i l}^{(j)} = \begin{cases}\frac{1}{n_{w_i}^{G(j)}(t)}, & \text { if }, j \in near by ~ l \\
% 0, & \text { if }, j \notin near by ~i
% \end{cases}  \\ 
% \text{self predict } w_{i i}^{(j)} = \begin{cases}\frac{1}{n_{w_i}^{G(j)}(t)}, & \text { if }, j \in near by ~ l \\
% 0, & \text { if }, j \notin near by ~i
% \end{cases}  \\
% \text{local correct } w_{i 0}^{(j)} = \begin{cases} t_y, & \text { if }, j \in near by ~ i \\
% 0, & \text { if }, j \notin near by ~i
% \end{cases}  
% \end{cases}
% \end{equation}

% ----------------------------

% Note that the dependence of the weights $ w_{i l}^{(j)}(t)$ on the trust observation history $\mathcal{O} (t)$ comes in through the choice of time dependent trusted neighborhood $N_i(t)$ [cf. (7)]. We use the parameter $\kappa$ later on, in our analysis of the dynamics observer in \eqref{dis_obs}, to obtain an upper bound on the entries of the matrix $W(t)$



% \begin{remark}
% The simulation will explore the use of two different design weights. In general, the design weight that prioritizes the local estimation term, as represented in equation \eqref{weight_0}, is used when a platoon wants to change lanes, necessitating an instantaneous adjustment over the consensus. In contrast, the design weight that distributed equally \eqref{weight_equal} is utilized for longitudinal maneuvers , that try to reach the consensus.    
% \end{remark}


\section{Estimation Error Dynamics}
In this part, we explore the conditions that guarantee the stability of the
estimation error dynamics of the proposed observer, as given in (5).

Using the error definitions:
\begin{align*}
e_0^{(j)}(k) = \bar{x}_i^{(j)}(k) - x_j(k),
\\
e_i^{(j)}(k) = \hat{x}_i^{(j)}(k) - x_j(k)
\end{align*}
The estimation error for the $j^{th}$ vehicle state as observed by vehicle $i$ is defined as:
\begin{align}
    e_i^{(j)}(k+1) = A\left[e_i^{(j)}(k) + \sum_{l \in \mathcal{N}_i} w_{il}^{(j)} (e_l^{(j)}(k) - e_i^{(j)}(k))\right. \notag \\
    +\left.  w_{i0}^{(j)}(e_0^{(j)}(k) - e_i^{(j)}(k))\right]. 
\end{align}
for the followers $( i, j \in V )$.


\textcolor{blue}{We can write this way to forget about the local error estimation, since local observer have good estimation,$e_0^{(j)}(k) = 0$ }

{\color{blue}
\begin{align}
    e_i^{(j)}(k+1) &= A\left[e_i^{(j)}(k) + \sum_{l \in \mathcal{N}_i} w_{il}^{(j)} (e_l^{(j)}(k) - e_i^{(j)}(k))\right. \notag \\
    +\left.  w_{i0}^{(j)} e_i^{(j)}(k))\right] \\ \notag 
    &= A\left[(1 - \sum_{l \in \mathcal{N}_i} w_{il}^{(j)} - w_{i0}^{(j)})e_i^{(j)}(k) + \sum_{l \in \mathcal{N}_i} w_{il}^{(j)} e_l^{(j)}(k)\right]\\ \notag
    &= A\left[ w_{ii}^{(j)}e_i^{(j)}(k) + \sum_{l \in \mathcal{N}_i} w_{il}^{(j)} e_l^{(j)}(k)\right]
\end{align}
--------------
}

Let $e^{(j)}(k)=\left[e_0^{(j)}(k)^{\top}, e_1^{(j)}(k)^{\top}, \ldots, e_m^{(j)}(k)^{\top}\right]^{\top}$. Based on the error dynamics (8), we obtain the following composite error dynamics:
$e_i^{(j)}(k+1) = R^{(j)} e^{(j)}(k) $

The matrix $R^{(j)}$ describes how the estimation errors evolve within the network. It has the following block structure:

$$
R^{(j)}=\left[\begin{array}{cccc}
R_{00}^{(j)} & R_{01}^{(j)} & \cdots & R_{0 m}^{(j)} \\
R_{10}^{(j)} & R_{11}^{(j)} & \cdots & R_{1 m}^{(j)} \\
\vdots & \vdots & \ddots & \vdots \\
R_{m 0}^{(j)} & R_{m 1}^{(j)} & \cdots & R_{m m}^{(j)}
\end{array}\right]
$$

where:
\begin{itemize}
    \item  $R_{00}^{(j)}=A-F_j C_{j, j}$ : This block captures the error dynamics of the virtual node (vehicle $j$ 's own estimate).
    \item $R_{0 i}^{(j)}=0$ for $i \neq 0$ : The virtual node's error does not directly depend on other vehicles' errors.
    \item $R_{i 0}^{(j)}=A w_{i 0}^{(j)}:$ This block captures the influence of the virtual node's error on vehicle $i$ 's error.
    \item Let $\left(1-\sum_{l \in N_i} w_{i l}^{(j)}-w_{i 0}^{(j)}\right) = w_{i i}^{(j)}$ 
    , so $R_{i i}^{(j)}=A w_{i i}^{(j)}$ : This block captures the contribution of vehicle $i$ 's own error. 
    \item $R_{i l}^{(j)}=A w_{i l}^{(j)}$ for $l \in N_i$ : This block captures the influence of neighboring vehicles' errors vehicle $i$ 's error.
\end{itemize}

Or we can write compact form : 
$$
\begin{aligned}
e^{(j)}(k+1)= & \underbrace{\left[\begin{array}{c|c}
A-F_j C_{j, j} & \mathbf{0} \\
\hline O^{(j)} \otimes A & Q^{(j)} \otimes A
\end{array}\right]}_{R^{(j)}} e^{(j)}(k) 
\end{aligned}
$$
The conditions on \( F_j \) and \( w_{il}^{(j)} \) ensure that the estimation error for all vehicles converges asymptotically to zero.

\textcolor{red}{\textbf{Question: Since the neighbor changes randomly, how many vehicles in the graph can be removed but the observer can still estimate the global state, and what is the consequence to the performance} }

\textcolor{red}{ can use switched graph to prove , but we don't know the  d time switch}
\textcolor{red}{....}
\textcolor{blue}{ I think it can not prove that, Impossible  !!!!! :(}

\subsection{Ideal stability analysis}
Notations : For a square matrix $M \in R^{n\times n}$,we use the symbols $sp(M)={\lambda \in C|det (\lambda I - M ) = 0}$ the set of all eigenvalues and $ \rho(M)$ to denote and the spectral radius of M the maximal magnitude of the eigenvalues of a real square matrix M, respectively.

Analyze about $Q^{j}$, everything is fine, no attack, so $ Q^{(j)}$ is Schur stable, indicated by a strongly connected intervehicle network according to Lemma2, combining the fact that $sp(A) \leq 1 $ for each $j \in V$, $Q(j) \otimes A$ is Schur stable. 
\\ 

\textbf{But} here we have attack, and the neighbor of $i$ changes, so the $ w_{il} , w_{i 0}$ change , lead to $Q^{(j)}$ change . So need to check the matrix $Q^{(j)}$ is Schur stable, i.e.,all the eigenvalues of $Q^{(j)}$ have modulus smaller than 1

 \subsection{First Attempt to prove Stability by the algebraic connectivity  }
Algebraic Connectivity and Its Importance in Distributed Estimation \cite{algebraic_connectivity}

The \textbf{algebraic connectivity} \( \lambda_2 \) is the second smallest eigenvalue of the \textbf{Laplacian matrix} \( L(G) \) of a graph \( G \). It characterizes how well the graph remains connected and directly influences:

\begin{enumerate}
    \item \textbf{Connectivity Condition}:
    \begin{itemize}
        \item \( \lambda_2 > 0 \) ensures that the communication graph is connected.
        \item If \( \lambda_2 = 0 \), the graph is disconnected, meaning there exist disjoint subgroups that do not exchange information.
    \end{itemize}
    
    \item \textbf{Consensus Rate}:
    \begin{itemize}
        \item The convergence rate of the consensus process depends on \( \lambda_2 \).
        \item A larger \( \lambda_2 \) results in \textbf{faster convergence}.
    \end{itemize}
    
    \item \textbf{Resilience to Switching Graphs}:
    \begin{itemize}
        \item If the graph is \textbf{switching} (i.e., the connectivity changes over time due to dynamic trust updates or malicious nodes), then we must analyze the \textbf{joint spectral properties} of the Laplacian matrices across different time steps.
    \end{itemize}
\end{enumerate}

{Switching Graphs and Convergence Conditions}

In a \textbf{switching graph}, the Laplacian matrix \( L(G) \) changes over time:

\[
L(t) = L(G_t)
\]

where \( G_t \) is the graph at time \( t \). Instead of a fixed \textbf{static} connectivity measure, we analyze the \textbf{joint connectivity} across multiple graphs.

\textbf{Case 1:} Graph is Connected at Each Time Step
If each graph \( G_t \) satisfies \( \lambda_2(G_t) > 0 \) at all times, then the consensus process converges.

\textbf{Case 2:} Graph Becomes Disconnected Occasionally
\begin{itemize}
    \item If some graphs are \textbf{temporarily disconnected} (i.e., \( \lambda_2(G_t) = 0 \)), the system might \textbf{slow down} or get stuck.
    \item However, if the graph is \textbf{connected on average} over time, then we can still guarantee \textbf{asymptotic convergence} using results from \textbf{jointly connected graph sequences}.
\end{itemize}

\textbf{Case 3}: Graph is Only Connected Over a Time Window
\begin{itemize}
    \item Instead of being \textbf{always connected}, the graph could satisfy the \textbf{joint connectivity condition}, meaning that over a sufficiently long window \( T \), the union of graphs \( G_t \) remains connected.
    \item \textbf{Key Result:} If there exists a time interval \( T \) such that the union of graphs \( G_{t}, G_{t+1}, ..., G_{t+T} \) has \( \lambda_2 > 0 \), then consensus is \textbf{still guaranteed}.
\end{itemize}

\textbf{Spectral Conditions for Convergence}

To guarantee \textbf{distributed estimation convergence}, we require that:

\[
\lim_{k \to \infty} W(k) W(k-1) ... W(0) x(0) = \mathbf{1} v^T x(0)
\]

where \( W(k) \) is the consensus matrix at time \( k \), and \textbf{all eigenvalues except one} of \( W(k) \) must remain within the unit circle.

The spectral radius condition ensures:

\[
\sup_k \rho(W(k) - \mathbf{1} v^T) < 1
\]

If switching occurs too frequently \textbf{without maintaining connectivity}, then \( \lambda_2 \to 0 \), slowing convergence and potentially \textbf{preventing consensus}.





The \textbf{speed of convergence} in discrete-time is determined by the second largest eigenvalue $\mu_2$ of $P$: \cite{Consensus_discret_Saber}

\begin{equation}
    \mu_2 = 1 - \varepsilon \lambda_2
\end{equation}

where $\lambda_2$ (the \textbf{algebraic connectivity}) is the second smallest eigenvalue of $L$. A larger $\lambda_2$ results in \textbf{faster} consensus.

For \textbf{continuous-time}, the convergence rate is:

\begin{equation}
    \| x(t) - x^* \| \leq e^{-\lambda_2 t} \| x(0) - x^* \|
\end{equation}

For \textbf{discrete-time}, the error decays as:

\begin{equation}
    \| x(k) - x^* \| \leq \mu_2^k \| x(0) - x^* \|
\end{equation}

which shows that \textbf{smaller $\mu_2$} leads to faster consensus.

By so we have this assumption : 
\begin{assumption}\label{number_attack}
Number attack vehicle is less equal ...
\end{assumption}
So by the assumption \ref{number_attack} the matrix $Q^{(j)}$ is Schur stable. So the the estimation error will converge.

% \begin{remark}
% % {Practical Considerations for Your System}
% \begin{itemize}
%     \item \textbf{Ensure Persistent Connectivity:} If trust values cause edges to drop dynamically, ensure that the remaining trusted agents form a connected graph \textbf{over time}.
%     \item \textbf{Minimize Large Connectivity Breaks:} Large gaps in communication where \( \lambda_2(G_t) = 0 \) should be avoided.
%     \item \textbf{Optimize Weights Based on Trust:} Weight adaptation must \textbf{not disconnect the network}, even when filtering malicious agents.
% \end{itemize}
% \end{remark}


% \subsection{Seconde Attempt to prove Stability by looking in the effects of malicious vehicles }


\section{After Check}
Since our calculation trust is not fully touched all vehicles, also the step aggerate trust using trust send over network. So we really not 100 $ \% $ sure about our global estimation. So we need a to check again our estimation output of distributed observer. We use pure predicted model with the previous value and neglect all the data recveide. 

\begin{equation}\label{dis_obs_predic}
\begin{aligned}
\hat{{z}}_i^{(j)}(k+1)= & A\left(\hat{x}_i^{(j)}(k)+w_{i 0}^{(j)}\left(\bar{x}_j(k)-\hat{x}_i^{(j)}(k)\right)\right)+B u_j(k), i, j \in V
\end{aligned}
\end{equation}

where $\bar{x}_j(k) = 0$ when $j \neq i$ 
\textbf{Consistency Error}: Calculate the error between the distributed observer and predict value 
    \begin{align}
    e_{i}^{(j)}(k) = \left| \hat{z}_i^{(j)}(k) - \hat{x}_i^{(j)}(k) \right| \end{align}
if  $e_{i}^{(j)}(k) \in Range$
The range could be large and cover all possibilities in the sense of physical.
\section{Simulation}



In this section, we validate the theoretical results and investigate the driving state estimation performance through extensive
simulations. Specifically, we apply our distributed driving state
observer \eqref{observer_dis_plug} to an open 2-nearest neighbor platoon. Moreover,
combining the proposed observer \eqref{observer_dis_plug} with the results of [15], we
propose an distributed observer-based controller and provide a
numerical simulation to illustrate its performance.

\subsection{Comparasion with \cite{Plug_play}}



\subsection{Senarior dedicated to the method}

\begin{figure}[htp]
    \centering
    \includegraphics[scale=0.4]{img/2Case_sim.png}
    \caption{Simulation case: "Left image" Perform a platoon try to match the speed and heading angle, "Right image" The platoon try to change the lane, vehicle following behavior.}
    \label{fig:2Case_sim}
\end{figure}

Platoon Control (Left Image):

Vehicles must maintain a safe following distance.
Speed synchronization is crucial to avoid abrupt braking or acceleration.
Heading angles should be matched to ensure stability.
Lane Change Coordination (Right Image):

The lead vehicle initiates the lane change, and the following vehicles must react accordingly.
Communication between vehicles (via V2V communication) is beneficial for smooth transitions.
The vehicle-following behavior ensures that the platoon does not break apart while changing lanes.




\paragraph{Expected Distance in Global Value}

The goal is to verify if the data aligns with expectations.

Its will look like \ref{local_sensor} , but instead compare with the sensor , we will compare with the expected value. 


 \textbf{Expected Consistency}: For neighbor \( l \)'s global estimate to be plausible, it should satisfy:

 
\begin{equation}
    \hat{X}_l^{(j)}(k) - \hat{X}_l^{(i)}(k) \approx d^{expect}_{i,j}(k)
\end{equation}
    

\textbf{Adaptive Expected Differences} 
The $d^{expect}_{i,j}$ here could be the distance and the velocity. 
\textcolor{blue}{So how to get that , we need to refer to the controller}
Since vehicles use a mix of IDM and CIDM, and you know which vehicles use which strategy (via 
$\gamma$), the expected differences (distance, velocity, acceleration) between consecutive vehicles should adapt to their control modes. Here's how to determine them:

\subsection*{1. Expected Distance (Spacing)}
The expected spacing depends on the following vehicle's control strategy, as it dictates the desired gap.

\begin{itemize}
    \item \textbf{IDM Steady-State Spacing}: When \(a_{\text{IDM}} = 0\) and \(\Delta v = 0\):

    \[ 1 - \left( \frac{v}{v_0} \right)^\delta = \left( \frac{s_0 + T v}{s} \right)^2 \]

\begin{align}
    s_{\text{IDM}} = \frac{s_0 + T v_0}{\sqrt{1 - \left( \frac{v}{v_0} \right)^\delta}}
\end{align}
    
    % If \(v = v_0\) and \(\delta = 4\):

    % \[ s_{\text{IDM}} = s_0 + T v_0 \]

    \item \textbf{CIDM Steady-State Spacing}: When \(a_{\text{CIDM}} = 0\) and \(\Delta \tilde{v}_{n-j+1} = 0\):

    \[ 1 - \left( \frac{v_n}{v_0} \right)^4 = \left( \frac{s_0 + T v_n}{\sum_{j=1}^N \alpha_{n j} w_{n j} \tilde{s}_{n-j+1}} \right)^2 \]

    % If \(v_n = v_0\):

    % \[ \sum_{j=1}^N \alpha_{n j} w_{n j} \tilde{s}_{n-j+1} = s_0 + T v_0 \]

\begin{align}
    s_{\text{CIDM}} = \frac{s_0 + T v_0}{\sqrt{1 - \left( \frac{v}{v_0} \right)^\delta}\cdot \sum_{j=1}^N j w_{n j}}
\end{align}

But we can simplified $s_{\text{IDM}}$ and $s_{\text{CIDM}}$ function by using a approximated as $s_0 + h_n v_0 $ where $h$ is a new time headway.

To make \( h \) concrete and adaptive, let's define it based on the vehicle's position in the platoon. The idea is that vehicles further back can maintain a smaller headway because they benefit from more information accumulated from vehicles ahead. Here's a simple, practical function:

\[ h_n = h_{\text{base}} - \Delta h \cdot (n - 2) \]

Where:
\begin{itemize}
    \item \( h_n \): Time headway for vehicle \( n \) (in seconds).
    \item \( h_{\text{base}} = T \): Base headway for the first follower (vehicle 2).
    \item \( \Delta h \): Headway reduction per position in the platoon.
    \item \( n \): Vehicle index (2, 3, 4 for followers; vehicle 1 is the leader and doesn't use \( h \)).
\end{itemize}

    \item \textbf{Mixed Control}: Since \(a_{\text{target}} = \gamma a_{\text{IDM}} + (1 - \gamma) a_{\text{CIDM}}\), the expected spacing is a blend:

    \[ s_{\text{expected}} = \gamma \cdot s_{\text{IDM}} + (1 - \gamma) \cdot s_{\text{CIDM}} \]

    Assuming \(s_{\text{IDM}} = s_0 + h_{base} v\) and \(s_{\text{CIDM}} = s_0 + h v\) :

    \[ s_{\text{expected}} = s_0 + [\gamma h_{base} + (1 - \gamma) h] v \]
\end{itemize}



\subsection*{2. Expected Velocity Difference}
In steady state, the platoon should have all vehicles moving at the same velocity (e.g., \(v_0\)), so:

\[ \text{expected\_diff\_velocity}_n = 0 \]

This holds regardless of IDM or CIDM, as both aim for velocity alignment in steady conditions.


\bibliographystyle{IEEEtran}


%=================
\bibliography{biblio}
%=================
\end{document}




% \end{document}
